\section{Related Work}
\label{sec:related}
% =============================================================

\subsection{Short-term solar irradiance forecasting}

Short-term GHI forecasting methods span a wide spectrum from statistical time-series
models to physically-based NWP and data-driven machine learning
\cite{diagne2013,voyant2017}. For intra-hour and sub-day horizons (0--6~h),
satellite-based and ground-observation methods generally outperform NWP, which
requires several hours of spin-up time \cite{perez2013}.

\textbf{Persistence} (assuming future irradiance equals the most recent observation)
is the standard benchmark for short-term forecasting. It is competitive at very
short horizons (under 30~minutes) when irradiance changes slowly, but degrades
rapidly beyond one hour \cite{lonij2013}.

\textbf{Statistical models} such as ARIMA and its seasonal variant SARIMA have
been applied to GHI series after removing the deterministic clear-sky envelope
\cite{reikard2009,bacher2009}. Working on the clear-sky index $k_t = \text{GHI}/\ghics$
improves stationarity and allows standard time-series machinery at hourly resolution.
Despite their simplicity, seasonal ARIMA models can be competitive at multi-hour
horizons in climates with regular diurnal patterns \cite{pedro2012}.

\subsection{Deep learning with satellite imagery}

CNN-based architectures applied to sequences of geostationary satellite patches
have become the dominant approach for satellite-based GHI forecasting. Early
work showed that cloud-motion vectors extracted from visible channels improve
short-term forecasts \cite{chow2011}. More recent studies feed raw satellite
reflectances directly into CNNs followed by fully-connected regressors or LSTMs
\cite{huertas2018,paletta2021,feng2022}. The encoder-decoder design, where a CNN
extracts per-frame spatial embeddings that are then processed by a recurrent model,
has shown consistent gains over purely spatial or purely temporal baselines
\cite{wang2019}.

Attention-based architectures including Vision Transformers (ViT) and Swin
Transformers have also been applied to satellite-based solar forecasting
\cite{acikgoz2022,liu2023}, achieving competitive performance at the cost of
increased parameter count and data requirements. For moderate dataset sizes such
as those encountered in single-site studies, CNN+LSTM hybrids remain a practical
and well-characterised baseline.

\subsection{Graph neural networks for spatial data}

GraphSAGE \cite{hamilton2017} is an inductive GNN that learns aggregation
functions over local neighbourhoods, enabling generalisation to unseen nodes.
It has been applied to traffic forecasting \cite{li2018}, air quality prediction,
and wind power generation \cite{kazem2021}, where spatial proximity defines the
graph topology. In the context of satellite imagery, representing pixels as graph
nodes allows the model to capture irregular spatial dependencies that periodic
CNN filters may miss, such as cloud boundaries or fractional cloud cover patterns.

Closest to our work, \cite{ji2022} applied a GNN-LSTM model to multi-site solar
irradiance forecasting, treating monitoring stations as nodes. Our approach differs
in that we build the graph within a single satellite patch (pixel level), without
requiring multi-station measurements, which is more practical for single-site
deployment.

\subsection{Forecasting in tropical and Andean climates}

Most GHI forecasting benchmarks are conducted in temperate or semi-arid climates
(Europe, southwestern United States, Australia) \cite{paletta2022,benchmarks2023}.
Tropical Andean sites present distinct challenges: convective cloud formation is
driven by topographic lifting and moisture advection, producing strong intraday
variability concentrated in the afternoon \cite{poveda2004}. Few published studies
address satellite-based forecasting for Colombian or equivalent tropical mountain
conditions, motivating the present work.

\subsection{Forecasting for power system operation and decision-making}

Accurate short-term irradiance forecasts feed directly into operational
decision-making for solar-integrated power systems: they support reserve
scheduling, unit commitment, and demand-response triggers, reducing the cost
of accommodating variable renewable generation \cite{energyrev2021}. For
renewable integration specifically, accurate forecasting reduces the need for
spinning reserves and allows pre-emptive load balancing, directly limiting the
impact of generation uncertainty on grid stability.

Beyond point forecasts, operators benefit from knowing \emph{why} a forecast
is uncertain. A forecast can be uncertain because the underlying cloud field
is intrinsically stochastic and would remain so even with infinite training
data (\emph{aleatoric} uncertainty), or because the model has seen too little
data resembling the current conditions and would become more confident with
more observations (\emph{epistemic} uncertainty) \cite{kendall2017}. The two
sources call for different operational responses --- accepting irreducible
variance during well-characterised but volatile conditions versus
provisioning extra reserve margin under genuinely novel, under-sampled
regimes --- a distinction that a single symmetric or distribution-free
interval cannot make on its own.

\subsection{Uncertainty quantification in deep learning for forecasting}

Bayesian deep learning provides a principled route to decomposing predictive
uncertainty. \cite{kendall2017} formalise the aleatoric/epistemic split for
deep regression and segmentation, combining a heteroscedastic Gaussian
negative-log-likelihood (NLL) head for aleatoric noise with Monte Carlo
sampling over model weights for epistemic uncertainty. Approximate posterior
sampling over weights can be obtained at training-time cost via Stochastic
Gradient Langevin Dynamics (SGLD, \cite{welling2011}), which injects
calibrated Gaussian noise into stochastic gradient updates so that, under a
decaying step size, the resulting trajectory asymptotically samples the true
posterior $p(\theta \mid \mathcal{D})$; collecting weight snapshots along a
single SGLD chain yields an inexpensive approximate ensemble whose spread
quantifies epistemic uncertainty, distinct from --- and complementary to ---
the aleatoric component captured by a heteroscedastic variance head. These
ideas have seen limited application
to satellite-based solar forecasting, where the dominant uncertainty driver
--- stochastic convective cloud formation --- is a textbook aleatoric source
that is rarely modelled explicitly alongside the model's own epistemic
limitations.

This work contributes to this literature by (i) developing spatially-aware
deep learning for solar forecasting at tropical Andean sites that remain
underrepresented in existing benchmarks, and (ii) applying an aleatoric/epistemic
uncertainty decomposition --- via a variance head and an SGLD posterior
ensemble --- to satellite-based GHI forecasting, supporting risk-stratified,
robust operational decisions that a single point forecast or symmetric
interval cannot provide.
